% §14 Exercises
%
% Cross-reference resolution status:
%
% DEFINED here:
%   \label{ex:ch14-1} through \label{ex:ch14-11}
%
% RESOLVED (used here, defined elsewhere):
%   \Cref{def:bhf}                                -- §14.2 (Plan 14 Task 3)
%   \Cref{sec:bhf-construction}                   -- §14.2 (Plan 14 Task 3)
%   \Cref{thm:bhf-optimality}                     -- §14.3 (Plan 14 Task 4)
%   \Cref{sec:space-analysis}                     -- §14.3 (Plan 14 Task 4)
%   \Cref{prop:bhf-fpr}                           -- §14.4 (Plan 14 Task 5)
%   \Cref{sec:fpr-analysis}                       -- §14.4 (Plan 14 Task 5)
%   \Cref{sec:construction-time}                  -- §14.5 (Plan 14 Task 6)
%   \Cref{def:bloom-filter}                       -- chapter 5 §5.2
%   \Cref{prop:bloom-overhead}                    -- chapter 8 §8.4
%   \Cref{thm:entropy-floor}                      -- chapter 8 §8.3
%   \Cref{ch:perfect-hash-filters}                -- chapter 15 stub
%
% Difficulty-marker macros (matching chapters 1-13):
\providecommand{\exstar}{$\star$}
\providecommand{\exstarstar}{$\star\star$}
\providecommand{\exstarstarstar}{$\star\star\star$}

\section*{Exercises}
\addcontentsline{toc}{section}{Exercises}

\begin{enumerate}

% -----------------------------------------------------------------------
% WARM-UP  (★ × 3): exercises 1, 2, 3
% -----------------------------------------------------------------------

\item \exstar
\label{ex:ch14-1}
\textbf{Fingerprint bits for target FPR.}
For each target false-positive rate $\fprate \in \{0.01, 0.001,
10^{-6}\}$, compute the BHF fingerprint width
$r = \lceil -\log_2 \fprate \rceil$. Verify that $2^{-r} \le
\fprate$ in each case, and report the over-achievement factor
$2^{-r} / \fprate$. The construction of \Cref{sec:bhf-construction}
predicts $r = 7, 10, 20$ for the three targets; report the
corresponding $2^{-r}$ values and confirm.

\item \exstar
\label{ex:ch14-2}
\textbf{BHF space at $n = 10^6, \fprate = 0.001$.}
Compute the BHF's fingerprint-table space in bits and in megabytes
for a designated set of size $n = 10^6$ at target $\fprate =
0.001$. (Ignore the sublinear perfect-hash overhead.) Compare to
the standard Bloom filter at the same $n$ and target, using the
formula $m = n \cdot \log_2(1/\fprate) \cdot \log_2 e$ from
\Cref{prop:bloom-overhead}. Report the absolute and relative
savings.

\item \exstar
\label{ex:ch14-3}
\textbf{Verify FNR = 0 in two sentences.}
Explain in at most two sentences why the BHF of \Cref{def:bhf}
cannot produce a false negative. The argument should refer to the
build pass that populates $F[\phi(a)] = g(a)$ for every $a \in A$
and to the fact that $\phi$ is a function, so $\phi(x)$ at query
time is the same slot $\phi(x)$ at build time.

% -----------------------------------------------------------------------
% MEDIUM  (★★ × 5): exercises 4, 5, 6, 7, 8
% -----------------------------------------------------------------------

\item \exstarstar
\label{ex:ch14-4}
\textbf{Notebook-tied BHF build.}
Use the chapter~14 Python notebook to build a BHF for $n = 10000$
items at target $\fprate = 0.01$, giving $r = 7$ fingerprint bits
per element.
\begin{enumerate}
  \item[(a)] Construct a designated set $A$ of $10000$ random
    32-bit integers and build the BHF using the notebook's
    perfect-hash construction.
  \item[(b)] Generate $10^5$ random query points drawn from outside
    $A$ (e.g., from a disjoint integer range) and count the number
    of queries that return $\top$.
  \item[(c)] Compute the empirical false-positive rate as
    (positives) / (total queries) and verify it lies within
    $\pm 0.001$ of the theoretical $2^{-7} \approx 0.00781$. State
    whether the agreement is consistent with the uniformity
    assumption of \Cref{prop:bhf-fpr}.
\end{enumerate}

\item \exstarstar
\label{ex:ch14-5}
\textbf{BHF vs Bloom space ratio.}
Derive the asymptotic ratio of BHF space to Bloom space at any
target $\fprate$, ignoring the sublinear perfect-hash overhead and
the integer-rounding gap.
\begin{enumerate}
  \item[(a)] Write the BHF bits-per-element as $r =
    \lceil -\log_2 \fprate \rceil \approx -\log_2 \fprate$ at small
    $\fprate$.
  \item[(b)] Write the Bloom bits-per-element as
    $\log_2(1/\fprate) \cdot \log_2 e$ from
    \Cref{prop:bloom-overhead}.
  \item[(c)] Compute the ratio (BHF) / (Bloom) and show it tends
    to $1 / \log_2 e \approx 0.693$ as $\fprate \to 0$. The BHF
    uses about $69.3\%$ of the bits the Bloom filter uses at the
    same target, the chapter's headline space saving.
\end{enumerate}

\item \exstarstar
\label{ex:ch14-6}
\textbf{Rounding gap.}
The integer rounding in $r = \lceil -\log_2 \fprate \rceil$
produces a worst case when $-\log_2 \fprate$ falls just above an
integer. Take $\fprate = 0.06$, so $-\log_2 0.06 \approx 4.06$ and
$r = 5$.
\begin{enumerate}
  \item[(a)] Compute the achieved FPR, $2^{-5} = 0.03125$, and the
    over-achievement factor $2^{-5} / 0.06 \approx 0.52$. The
    BHF over-achieves the target by approximately a factor of two
    in this case.
  \item[(b)] Quantify the wasted bits: the entropy floor is $4.06$
    bits per element, the BHF uses $5$ bits per element, so $0.94$
    bits per element are wasted to integer rounding. At $n = 10^6$,
    this is $0.94 \cdot 10^6 = 940{,}000$ bits or roughly $115$
    kilobytes of waste.
  \item[(c)] Compare this to the rounding waste at $\fprate =
    2^{-5} = 0.03125$ (where the target lands exactly on an integer
    bit boundary). State which target the designer should choose
    when the precise target is negotiable; the chapter's
    recommendation in \Cref{sec:fpr-analysis} is to round the
    target rate down to the nearest power of two when possible.
\end{enumerate}

\item \exstarstar
\label{ex:ch14-7}
\textbf{Perfect-hash overhead.}
The space analysis of \Cref{sec:space-analysis} cited the
Belazzougui-Botelho-Dietzfelbinger (BDZ, 2009) bound
$L_\phi = O(n \log \log \log n / \log n)$ for the perfect-hash
representation.
\begin{enumerate}
  \item[(a)] At $n = 10^6$, evaluate the formula. Take $\log n = 20$
    in base $2$, $\log \log n \approx 4.32$, $\log \log \log n
    \approx 2.11$. Compute the per-element overhead $L_\phi / n
    \approx 2.11 / 20 \approx 0.1$ bits per element in the
    asymptotic form. Practical BDZ implementations report between
    $2$ and $3$ bits per element at $n = 10^6$ owing to the hidden
    asymptotic constants. Report both the formula value and the
    practical value.
  \item[(b)] At $n = 10^9$, evaluate the same formula. Take $\log n
    \approx 30$, $\log \log n \approx 4.91$, $\log \log \log n
    \approx 2.29$. Compute $L_\phi / n \approx 2.29 / 30 \approx
    0.076$ bits per element in the asymptotic form. Practical
    implementations report between $1.5$ and $2.5$ bits per
    element.
  \item[(c)] State the per-element overhead's behavior as $n \to
    \infty$ and explain why the BHF's total per-element cost is
    asymptotically $r$ bits, the fingerprint width, with the
    perfect-hash term vanishing.
\end{enumerate}

\item \exstarstar
\label{ex:ch14-8}
\textbf{Composition under BHF.}
Two Bloom filters on disjoint sets $A_1$ and $A_2$ with the same
parameters can be merged into a Bloom filter on $A_1 \cup A_2$ by
bitwise OR of the bit arrays. Two BHFs on disjoint sets cannot be
merged this way. Explain why.
\begin{enumerate}
  \item[(a)] Identify the structural reason: the perfect hash
    $\phi_1$ for $A_1$ is a bijection $A_1 \to [|A_1|]$, and the
    perfect hash $\phi_2$ for $A_2$ is a bijection $A_2 \to
    [|A_2|]$. The two address spaces overlap on the integers
    $\{1, \ldots, \min(|A_1|, |A_2|)\}$ but have no shared
    semantics; slot $i$ in $\phi_1$'s table stores the fingerprint
    of some $a_1 \in A_1$, while slot $i$ in $\phi_2$'s table
    stores the fingerprint of a different element $a_2 \in A_2$.
    A bitwise OR (or any pointwise combination) of the two
    fingerprint tables produces an artifact that does not
    correspond to any valid BHF on $A_1 \cup A_2$.
  \item[(b)] Sketch the BHF's merge procedure: rebuild a new
    perfect hash $\phi'$ on $A_1 \cup A_2$ and refill the
    fingerprint table from scratch. The cost is $O(|A_1| +
    |A_2|)$, the same as building a BHF on the merged set from
    raw input.
  \item[(c)] State the asymmetry's source: the Bloom filter's
    address space is data-independent (the bit array is a
    universal address space shared across all instances), so bit-OR
    merges are valid. The BHF's address space is data-dependent
    (the perfect hash is built from the designated set), so merges
    require rebuilding the address space. The data-dependent
    addressing is what gives the BHF its space optimality and what
    prevents the cheap merge.
\end{enumerate}

% -----------------------------------------------------------------------
% CHALLENGING / EXTENSION  (★★★ × 3): exercises 9, 10, 11
% -----------------------------------------------------------------------

\item \exstarstarstar
\label{ex:ch14-9}
\textbf{BHF FPR exact analysis.}
The proposition \Cref{prop:bhf-fpr} bounds the BHF's false-positive
rate by $2^{-r}$. Derive the exact FPR for a non-adversarial query
$x$ drawn uniformly at random from $\Universe \setminus A$, taking
into account the perfect-hash range-to-universe ratio.
\begin{enumerate}
  \item[(a)] Suppose the perfect hash $\phi$ is extended to
    $\Universe$ by a uniform hash to $[n]$ on inputs outside $A$.
    For a random $x \in \Universe \setminus A$, the probability
    that $\phi(x) \in \phi(A) = [n]$ is exactly $1$ in this
    extension (every slot is reached). Write the exact FPR as
    $\fprate_{\text{BHF}} = 1 \cdot 2^{-r} = 2^{-r}$, recovering
    the proposition.
  \item[(b)] Alternatively suppose $\phi$ is extended by mapping
    inputs outside $A$ to a sentinel slot $0 \notin [n]$, and the
    query returns $\bot$ on the sentinel. For a random $x \in
    \Universe \setminus A$ with universe size $|\Universe|$, the
    probability that $\phi(x) \neq 0$ depends on the extension
    rule's design. Suppose the extension routes a fraction
    $\alpha \in [0, 1]$ of non-$A$ inputs to a non-sentinel slot;
    write the exact FPR as $\fprate_{\text{BHF}} = \alpha \cdot
    2^{-r}$.
  \item[(c)] In a practical deployment with $|\Universe| \gg n$
    and a uniform-extension $\phi$, the value of $\alpha$ is close
    to $n / |\Universe|$ if the perfect-hash construction is
    well-extended. Show that for $|\Universe| = 2^{64}$ and $n =
    10^6$, $\alpha \approx 10^6 / 2^{64} \approx 5 \times 10^{-14}$,
    making the exact FPR vanishingly small in the sentinel-aware
    extension. Discuss why the proposition's upper bound $2^{-r}$
    is the right design target despite the achievable FPR being
    much smaller in practice: the upper bound is the adversarial
    worst case and is what the designer must budget against.
\end{enumerate}

\item \exstarstarstar
\label{ex:ch14-10}
\textbf{Dynamic BHF.}
The BHF of \Cref{def:bhf} is static: insertions require an $O(n)$
rebuild. Sketch a dynamic variant that supports amortized cheaper
insertions.
\begin{enumerate}
  \item[(a)] \emph{Tiered layout.} Maintain a static BHF on a
    ``base'' set $A_0$ and a small auxiliary Bloom filter on a
    ``delta'' set $\Delta$ of recent insertions. A query checks
    both: $\top$ iff the BHF returns $\top$ or the Bloom filter
    returns $\top$. The structure supports streaming insertions
    into the Bloom-filter delta in $O(k)$ time and queries in
    $O(1)$ time (BHF) plus $O(k)$ time (Bloom delta).
  \item[(b)] \emph{Rebuild trigger.} When $|\Delta|$ exceeds a
    threshold (say, $|\Delta| \ge 0.1 \cdot |A_0|$), rebuild the
    BHF on $A_0 \cup \Delta$ and reset the delta to empty. The
    rebuild costs $O(|A_0| + |\Delta|)$, amortized over the
    inter-rebuild insertion volume. Derive the amortized cost per
    insertion as $O(|A_0| / 0.1 |A_0|) = O(10) = O(1)$ in the
    base case, plus the per-insertion Bloom cost $O(k)$. The
    dynamic variant achieves amortized $O(k)$ insertion time.
  \item[(c)] \emph{Space cost.} The dynamic variant pays an extra
    $O(0.1 n \cdot \log_2(1/\fprate) \cdot \log_2 e)$ bits for the
    Bloom delta, a $14\%$ space overhead over the pure BHF at the
    $0.1$ rebuild threshold. Discuss the design trade-off between
    rebuild frequency and space overhead: smaller thresholds give
    smaller deltas (less space) at the cost of more frequent
    rebuilds (more amortized work per insertion).
\end{enumerate}

\item \exstarstarstar
\label{ex:ch14-11}
\textbf{Entropy-floor matching.}
Prove that any FNR-zero approximate set using $r$ bits per
fingerprint in a perfect-hash-addressed table achieves FPR
$\le 2^{-r}$ at best, establishing that the BHF is asymptotically
optimal among constructions of this form.
\begin{enumerate}
  \item[(a)] \emph{Setup.} Consider any FNR-zero data structure of
    the form $(\phi, F)$ where $\phi : A \to [n]$ is a perfect
    hash and $F : [n] \to \{0, 1\}^r$ is a fingerprint table. The
    FNR-zero requirement forces $F[\phi(a)]$ to deterministically
    indicate $a$'s membership for every $a \in A$; the only
    freedom is in how the fingerprint $F[\phi(a)]$ is computed.
  \item[(b)] \emph{The bound.} For a query $x \notin A$ mapped by
    $\phi$ to a slot $i \in [n]$, the table value $F[i]$ takes one
    of $2^r$ possible values, and the query's verdict is $\top$ iff
    the value matches some $r$-bit string the query algorithm
    derives from $x$. Under any non-adversarial query distribution
    and any uniformly-distributed fingerprint-derivation function,
    the match probability is $2^{-r}$. The bound is tight: the
    BHF achieves it; any construction achieving less must violate
    the uniformity of the fingerprint derivation or restrict the
    query distribution.
  \item[(c)] \emph{Combine with the entropy floor.} The bound $r
    \ge -\log_2 \fprate$ from \Cref{thm:entropy-floor} and the
    bound (FPR $\ge 2^{-r}$) from the present argument together
    pin the achievable $(r, \fprate)$ region for FNR-zero
    perfect-hash-addressed structures: $\fprate \ge 2^{-r}$ and
    $r \ge \lceil -\log_2 \fprate \rceil$. The BHF lives at the
    intersection: it uses exactly $r = \lceil -\log_2 \fprate
    \rceil$ bits per fingerprint and achieves exactly the FPR
    bound $2^{-r} \le \fprate$. State the optimality conclusion:
    no FNR-zero perfect-hash-addressed structure can use fewer
    bits per element than the BHF to meet the same FPR target,
    asymptotically in $n$.
\end{enumerate}

\end{enumerate}
